% ==============================================================================
% MODELO ACADÊMICO MODULAR EM LATEX — EDIÇÃO FACILITADA
% ==============================================================================
% Bem-vindo(a)! Este arquivo foi desenhado para que você escreva pré-projetos,
% projetos de pesquisa, TCCs, monografias e relatórios acadêmicos sem precisar
% dominar programação em LaTeX.
%
% Tudo o que está entre [COLCHETES] é um espaço reservado para o seu texto.
% Linhas começando com % são comentários explicativos (não saem no PDF).
%
% ------------------------------------------------------------------------------
% PASSO A PASSO PARA USAR ESTE MODELO
% ------------------------------------------------------------------------------
% 1. Na PARTE 1 (logo abaixo), preencha os dados do seu trabalho entre { }.
% 2. Na PARTE 2, escolha o que quer mostrar ou esconder usando {sim} ou {nao}.
% 3. Pule o "Motor de Formatação" e vá direto para a PARTE 3 e PARTE 4 para
%    escrever seus textos, adicionar, mover ou remover seções.
% 4. No final do arquivo (PARTE 6), há uma "Caixa de Ferramentas" com blocos
%    prontos para copiar e colar (Figuras, Tabelas ABNT, Quadros, Equações,
%    Citações longas, Listas, Modelos de Referências, etc.).
% 5. Como gerar o PDF:
%    - No computador (recomendado): tectonic modelo.tex
%    - No Overleaf: Menu (canto superior esquerdo) > Compilador > XeLaTeX
%
% ------------------------------------------------------------------------------
% GUIA RÁPIDO DE ESCRITA NO LATEX
% ------------------------------------------------------------------------------
% - NOVO PARÁGRAFO: Deixe uma linha totalmente em branco entre um parágrafo e
%   outro. Nunca use \\ para separar parágrafos comuns.
% - NEGRITO E ITÁLICO:
%     \textbf{texto em negrito}
%     \textit{texto em itálico}   (use para palavras estrangeiras e espécies)
% - NOTA DE RODAPÉ:
%     Texto normal\footnote{Texto da nota de rodapé.}
% - SÍMBOLOS ESPECIAIS (precisam de barra invertida antes se usados no texto):
%     \%   \&   \$   \#   \_   \{   \}
% - GRAUS E UNIDADES:
%     20 \textdegree{}C    ou    70 \(\pm\) 2 \textdegree{}C
% - SUBSCRITO E SOBRESCRITO QUÍMICO/MATEMÁTICO:
%     \(\mathrm{EC}_{50}\), \(\mathrm{H_2O}\), \(10^{-3}\)
% - DICA DE OURO: Nunca apague uma chave { ou } sozinha; elas sempre trabalham
%   em pares {assim}. Para achar o que falta preencher, busque por [ no editor.
% ==============================================================================

\documentclass[12pt,a4paper,oneside]{article}

% --- Inicialização interna dos comandos de configuração (não altere) ---
\makeatletter
\newcommand{\@instituicao}{}\newcommand{\Instituicao}[1]{\renewcommand{\@instituicao}{#1}}
\newcommand{\@unidade}{}\newcommand{\Unidade}[1]{\renewcommand{\@unidade}{#1}}
\newcommand{\@curso}{}\newcommand{\Curso}[1]{\renewcommand{\@curso}{#1}}
\newcommand{\@autor}{}\newcommand{\Autor}[1]{\renewcommand{\@autor}{#1}}
\newcommand{\@titulo}{}\newcommand{\Titulo}[1]{\renewcommand{\@titulo}{#1}}
\newcommand{\@subtitulo}{}\newcommand{\Subtitulo}[1]{\renewcommand{\@subtitulo}{#1}}
\newcommand{\@cidade}{}\newcommand{\Cidade}[1]{\renewcommand{\@cidade}{#1}}
\newcommand{\@ano}{}\newcommand{\Ano}[1]{\renewcommand{\@ano}{#1}}
\newcommand{\@natureza}{}\newcommand{\Natureza}[1]{\renewcommand{\@natureza}{#1}}
\newcommand{\@orientacao}{}\newcommand{\Orientacao}[1]{\renewcommand{\@orientacao}{#1}}
\newcommand{\@coorientacao}{}\newcommand{\Coorientacao}[1]{\renewcommand{\@coorientacao}{#1}}
\newcommand{\@arquivologo}{}\newcommand{\ArquivoLogo}[1]{\renewcommand{\@arquivologo}{#1}}
\newcommand{\@alturalogo}{2.7cm}\newcommand{\AlturaLogo}[1]{\renewcommand{\@alturalogo}{#1}}

\newcommand{\@definirchave}[2]{%
  \expandafter\newif\csname if#1\endcsname
  \csname #1#2\endcsname
  \expandafter\newcommand\csname #1\endcsname[1]{%
    \edef\@tempa{\detokenize{##1}}%
    \edef\@sims{\detokenize{sim}}\edef\@simS{\detokenize{SIM}}\edef\@simCap{\detokenize{Sim}}%
    \edef\@s{\detokenize{s}}\edef\@S{\detokenize{S}}\edef\@true{\detokenize{true}}%
    \ifx\@tempa\@sims\csname #1true\endcsname\else
    \ifx\@tempa\@simS\csname #1true\endcsname\else
    \ifx\@tempa\@simCap\csname #1true\endcsname\else
    \ifx\@tempa\@s\csname #1true\endcsname\else
    \ifx\@tempa\@S\csname #1true\endcsname\else
    \ifx\@tempa\@true\csname #1true\endcsname\else
      \csname #1false\endcsname
    \fi\fi\fi\fi\fi\fi
  }%
}
\@definirchave{MostrarCapa}{true}
\@definirchave{MostrarFolhaRosto}{true}
\@definirchave{MostrarDedicatoria}{false}
\@definirchave{MostrarAgradecimentos}{false}
\@definirchave{MostrarEpigrafe}{false}
\@definirchave{MostrarResumo}{true}
\@definirchave{MostrarAbstract}{false}
\@definirchave{MostrarListaFiguras}{false}
\@definirchave{MostrarListaQuadros}{false}
\@definirchave{MostrarListaTabelas}{false}
\@definirchave{MostrarListaSiglas}{false}
\@definirchave{MostrarSumario}{true}
\@definirchave{MostrarCronograma}{true}
\@definirchave{MostrarReferencias}{true}
\@definirchave{MostrarApendices}{false}
\@definirchave{MostrarAnexos}{false}
\@definirchave{SecoesEmNovaPagina}{true}
\@definirchave{EspacamentoUmEMeio}{true}
\@definirchave{ReiniciarPaginaNoCorpo}{false}
\makeatother


% ==============================================================================
% PARTE 1 — DADOS DO TRABALHO (EDITE APENAS O QUE ESTÁ ENTRE CHAVES { })
% ==============================================================================
% Estes dados preenchem automaticamente a Capa, a Folha de Rosto e os metadados
% do arquivo PDF. Se um campo opcional não se aplicar ao seu trabalho, basta
% deixá-lo vazio: {} (ele não deixará buracos nem linhas em branco na página).

\Instituicao{[NOME DA INSTITUIÇÃO]}
\Unidade{[UNIDADE, CENTRO OU DEPARTAMENTO]}      % Opcional: deixe {} para ocultar
\Curso{[NOME DO CURSO]}                          % Opcional: deixe {} para ocultar
\Autor{[NOME COMPLETO DO AUTOR]}
\Titulo{[TÍTULO DO TRABALHO]}
\Subtitulo{}                                     % Opcional: se preenchido, aparece após ": "
\Cidade{[CIDADE]}
\Ano{[ANO]}

% Descrição da natureza do trabalho (aparece recuada na Folha de Rosto):
\Natureza{[DESCRIÇÃO DO TRABALHO, FINALIDADE, CURSO E INSTITUIÇÃO]}

% Orientação e coorientação (inclua o rótulo desejado, ex.: Orientador: Prof. Dr. ...):
\Orientacao{[ORIENTAÇÃO: NOME E TITULAÇÃO]}
\Coorientacao{}                                  % Opcional: deixe {} para ocultar

% Logotipo ou brasão da instituição no topo da capa (opcional):
% Para usar, coloque o caminho da imagem entre chaves, ex.: {imagens/brasao.png}
% Deixe vazio {} para não exibir imagem.
\ArquivoLogo{}
\AlturaLogo{2.7cm}                               % Altura da imagem do logotipo


% ==============================================================================
% PARTE 2 — PAINEL DE CONTROLE: LIGAR OU DESLIGAR PARTES ({sim} ou {nao})
% ==============================================================================
% Para ativar uma página ou seção, escreva {sim}. Para ocultá-la sem precisar
% apagar seus textos, escreva {nao}.

% --- Páginas Iniciais (Pré-Textuais) ---
\MostrarCapa{sim}               % Capa institucional
\MostrarFolhaRosto{sim}         % Folha de rosto com natureza do trabalho e orientação
\MostrarDedicatoria{nao}        % Página de dedicatória (opcional)
\MostrarAgradecimentos{nao}     % Página de agradecimentos (opcional)
\MostrarEpigrafe{nao}           % Página de epígrafe / citação inicial (opcional)
\MostrarResumo{sim}             % Resumo em português
\MostrarAbstract{nao}           % Abstract / resumo em língua estrangeira (opcional)
\MostrarListaFiguras{nao}       % Lista automática de figuras (opcional)
\MostrarListaQuadros{nao}       % Lista automática de quadros (opcional)
\MostrarListaTabelas{nao}       % Lista automática de tabelas (opcional)
\MostrarListaSiglas{nao}        % Lista de abreviaturas e siglas (opcional)
\MostrarSumario{sim}            % Sumário automático

% --- Seções Especiais e Pós-Textuais ---
\MostrarCronograma{sim}         % Seção de cronograma de atividades
\MostrarReferencias{sim}        % Seção de referências bibliográficas
\MostrarApendices{nao}          % Apêndices (materiais elaborados pelo próprio autor)
\MostrarAnexos{nao}             % Anexos (documentos/materiais elaborados por terceiros)

% --- Opções Gerais de Formatação ---
\SecoesEmNovaPagina{sim}        % {sim} = cada seção principal começa em nova página
                                % {nao} = seções seguem na mesma página se houver espaço
\EspacamentoUmEMeio{sim}        % {sim} = espaçamento 1,5 entre linhas (padrão acadêmico)
                                % {nao} = espaçamento simples em todo o texto
\ReiniciarPaginaNoCorpo{nao}    % {nao} = conta desde a folha de rosto (padrão ABNT)
                                % {sim} = reinicia a numeração em 1 na primeira seção


% ==============================================================================
% MOTOR DE FORMATAÇÃO AUTOMÁTICA (NÃO É NECESSÁRIO EDITAR ESTA PARTE)
% ==============================================================================
% Pule diretamente para a PARTE 3 (após \begin{document}) para editar os textos.
% ------------------------------------------------------------------------------
\usepackage{fontspec}
\IfFontExistsTF{Times New Roman}
  {\setmainfont{Times New Roman}}
  {\setmainfont{TeX Gyre Termes}}
\usepackage[brazil]{babel}
\usepackage[left=3cm,right=2cm,top=3cm,bottom=2cm,headheight=14pt,headsep=18pt]{geometry}
\usepackage{setspace,graphicx,amsmath,array,titlesec,fancyhdr,enumitem}
\usepackage[titles]{tocloft}
\usepackage{xurl,tabularx,unicode-math}
\setmathfont{texgyretermes-math.otf}
\usepackage[hidelinks,unicode,hypertexnames=false,pdflang={pt-BR}]{hyperref}

% Tipos de colunas simplificados para Tabelas e Quadros:
% L = flexível à esquerda | C = flexível centralizada | R = flexível à direita
% P{largura} = largura fixa à esquerda | M{largura} = largura fixa centralizada
\newcolumntype{L}{>{\raggedright\arraybackslash}X}
\newcolumntype{C}{>{\centering\arraybackslash}X}
\newcolumntype{R}{>{\raggedleft\arraybackslash}X}
\newcolumntype{P}[1]{>{\raggedright\arraybackslash}p{#1}}
\newcolumntype{M}[1]{>{\centering\arraybackslash}m{#1}}

\makeatletter
\AtBeginDocument{%
  \hypersetup{pdftitle={\@titulo},pdfauthor={\@autor}}%
}
\urlstyle{same}
\addtolength{\textheight}{-5pt}
\ifEspacamentoUmEMeio\onehalfspacing\else\singlespacing\fi
\setlength{\parindent}{1.25cm}
\setlength{\parskip}{0pt}
\setlength{\emergencystretch}{2em}
\widowpenalty=10000
\clubpenalty=10000

% Cabeçalho e paginação
\fancyhf{}
\fancyhead[R]{\fontsize{10}{12}\selectfont\thepage}
\renewcommand{\headrulewidth}{0pt}
\pagestyle{empty}

% Formatação de títulos de seção (padrão acadêmico/ABNT)
\titleformat{\section}{\normalfont\bfseries\normalsize}{\thesection}{.5em}{}
\titleformat{\subsection}{\normalfont\bfseries\normalsize}{\thesubsection}{.5em}{}
\titleformat{\subsubsection}{\normalfont\normalsize}{\thesubsubsection}{.5em}{}
\titleformat{name=\section,numberless}{\normalfont\bfseries\normalsize\centering}{}{0pt}{}
\titlespacing{\section}{0pt}{0pt}{\baselineskip}
\titlespacing{\subsection}{0pt}{\baselineskip}{\baselineskip}
\titlespacing{\subsubsection}{0pt}{\baselineskip}{\baselineskip}

% Configuração do Sumário, Lista de Figuras, Lista de Tabelas e Lista de Quadros
\setcounter{tocdepth}{3}
\addto\captionsbrazil{%
  \renewcommand{\contentsname}{SUMÁRIO}%
  \renewcommand{\listfigurename}{LISTA DE FIGURAS}%
  \renewcommand{\listtablename}{LISTA DE TABELAS}%
  \renewcommand{\figurename}{Figura}%
  \renewcommand{\tablename}{Tabela}%
}
\renewcommand{\cfttoctitlefont}{\hfill\normalsize\bfseries}
\renewcommand{\cftaftertoctitle}{\hfill\mbox{}}
\renewcommand{\cftloftitlefont}{\hfill\normalsize\bfseries}
\renewcommand{\cftafterloftitle}{\hfill\mbox{}}
\renewcommand{\cftlottitlefont}{\hfill\normalsize\bfseries}
\renewcommand{\cftafterlottitle}{\hfill\mbox{}}

\renewcommand{\cftsecfont}{\bfseries}
\renewcommand{\cftsecpagefont}{\normalfont}
\renewcommand{\cftsecleader}{\normalfont\cftdotfill{\cftdotsep}}
\renewcommand{\cftsubsecfont}{\normalfont\bfseries}
\renewcommand{\cftsubsubsecfont}{\normalfont}
\renewcommand{\cftsubsecpagefont}{\normalfont}
\renewcommand{\cftsubsubsecpagefont}{\normalfont}
\setlength{\cftbeforesecskip}{3pt}
\setlength{\cftsecindent}{0pt}
\setlength{\cftsubsecindent}{0pt}
\setlength{\cftsubsubsecindent}{0pt}
\setlength{\cftsecnumwidth}{4em}
\setlength{\cftsubsecnumwidth}{4em}
\setlength{\cftsubsubsecnumwidth}{4em}

\setlength{\cftfigindent}{0pt}
\setlength{\cftfignumwidth}{5.5em}
\renewcommand{\cftfigpresnum}{Figura }
\renewcommand{\cftfigaftersnum}{ \textendash{} }

\setlength{\cfttabindent}{0pt}
\setlength{\cfttabnumwidth}{5.5em}
\renewcommand{\cfttabpresnum}{Tabela }
\renewcommand{\cfttabaftersnum}{ \textendash{} }

% Criação da Lista de Quadros automática
\newcommand{\listquadroname}{LISTA DE QUADROS}
\newlistof{quadro}{qua}{\listquadroname}
\renewcommand{\cftquatitlefont}{\hfill\normalsize\bfseries}
\renewcommand{\cftafterquatitle}{\hfill\mbox{}}
\setlength{\cftquadroindent}{0pt}
\setlength{\cftquadronumwidth}{5.5em}
\renewcommand{\cftquadropresnum}{Quadro }
\renewcommand{\cftquadroaftersnum}{ \textendash{} }

% Comandos auxiliares internos
\newcommand{\linhaopcional}[1]{%
  \if\relax\detokenize\expandafter{#1}\relax\else #1\par\fi
}
\newcommand{\@imprimirtitulocompleto}{%
  \textbf{\@titulo}%
  \if\relax\detokenize\expandafter{\@subtitulo}\relax\else
    : \normalfont\@subtitulo
  \fi
}

% Geração automática da Capa e da Folha de Rosto
\newcommand{\GerarCapaEFolhaDeRosto}{%
  \ifMostrarCapa
    \thispagestyle{empty}
    \begin{center}
      \if\relax\detokenize\expandafter{\@arquivologo}\relax\else
        \IfFileExists{\@arquivologo}{%
          \includegraphics[height=\@alturalogo]{\@arquivologo}\par\vspace{8pt}%
        }{%
          \fbox{\parbox{0.65\linewidth}{\centering\footnotesize\singlespacing
            [Logotipo não encontrado em \textbf{\@arquivologo}.\\
            Verifique o caminho em \textbf{\textbackslash ArquivoLogo} ou deixe \textbf{\{\}} vazio.]}}\par\vspace{8pt}%
        }%
      \fi
      \linhaopcional{\@instituicao}
      \linhaopcional{\@unidade}
      \linhaopcional{\@curso}
      \vspace{1.8cm}
      \@autor\par
      \vfill
      \@imprimirtitulocompleto\par
      \vfill
      \@cidade\par\@ano
    \end{center}
    \clearpage
  \fi
  % A capa não entra na contagem; a primeira página após a capa inicia em 1.
  \setcounter{page}{1}%
  \ifMostrarFolhaRosto
    \thispagestyle{empty}
    \begin{center}
      \@autor\par
      \vspace{3cm}
      \@imprimirtitulocompleto\par
    \end{center}
    \vspace{2cm}
    \begin{flushright}
      \begin{minipage}{.5\textwidth}
        \singlespacing\sloppy\hbadness=10000
        \@natureza\par
        \vspace{12pt}
        \linhaopcional{\@orientacao}
        \linhaopcional{\@coorientacao}
      \end{minipage}
    \end{flushright}
    \vfill
    \begin{center}\@cidade\par\@ano\end{center}
    \clearpage
  \fi
}

% Geração automática das Listas e do Sumário
\newcommand{\GerarListasESumario}{%
  \ifMostrarListaFiguras
    \thispagestyle{empty}
    \listoffigures
    \thispagestyle{empty}
    \clearpage
  \fi
  \ifMostrarListaQuadros
    \thispagestyle{empty}
    \listofquadro
    \thispagestyle{empty}
    \clearpage
  \fi
  \ifMostrarListaTabelas
    \thispagestyle{empty}
    \listoftables
    \thispagestyle{empty}
    \clearpage
  \fi
  \ifMostrarSumario
    \thispagestyle{empty}
    \tableofcontents
    \thispagestyle{empty}
    \clearpage
  \fi
  % Inicia o corpo do trabalho com numeração visível no topo direito
  \pagestyle{fancy}%
  \ifReiniciarPaginaNoCorpo
    \setcounter{page}{1}%
  \fi
}
\makeatother

% Ambientes para Páginas Pré-Textuais (ocultam-se sozinhos se marcados como {nao})
\NewDocumentEnvironment{BlocoDedicatoria}{+b}{%
  \ifMostrarDedicatoria
    \thispagestyle{empty}
    \vspace*{\fill}
    \begin{flushright}
      \begin{minipage}{0.5\textwidth}
        \singlespacing\sloppy\hbadness=10000\noindent #1
      \end{minipage}
    \end{flushright}
    \vspace{2cm}
    \clearpage
  \fi
}{}

\NewDocumentEnvironment{BlocoAgradecimentos}{+b}{%
  \ifMostrarAgradecimentos
    \thispagestyle{empty}
    \section*{AGRADECIMENTOS}
    #1
    \clearpage
  \fi
}{}

\NewDocumentEnvironment{BlocoEpigrafe}{m +b}{%
  \ifMostrarEpigrafe
    \thispagestyle{empty}
    \vspace*{\fill}
    \begin{flushright}
      \begin{minipage}{0.55\textwidth}
        \singlespacing\sloppy\hbadness=10000\noindent\textit{#2}\par\vspace{6pt}
        \hfill (#1)
      \end{minipage}
    \end{flushright}
    \vspace{2cm}
    \clearpage
  \fi
}{}

\NewDocumentEnvironment{BlocoResumo}{m +b}{%
  \ifMostrarResumo
    \thispagestyle{empty}
    \section*{RESUMO}
    \begingroup
    \singlespacing
    \noindent #2\par
    \vspace{\baselineskip}
    \noindent\textbf{Palavras-chave:} #1\par
    \endgroup
    \clearpage
  \fi
}{}

\NewDocumentEnvironment{BlocoAbstract}{m +b}{%
  \ifMostrarAbstract
    \thispagestyle{empty}
    \section*{ABSTRACT}
    \begingroup
    \singlespacing
    \noindent #2\par
    \vspace{\baselineskip}
    \noindent\textbf{Keywords:} #1\par
    \endgroup
    \clearpage
  \fi
}{}

\NewDocumentEnvironment{BlocoSiglas}{+b}{%
  \ifMostrarListaSiglas
    \thispagestyle{empty}
    \section*{LISTA DE ABREVIATURAS E SIGLAS}
    \begingroup
    \singlespacing
    \noindent
    \begin{tabularx}{\linewidth}{@{}P{3cm}L@{}}
      #1
    \end{tabularx}
    \endgroup
    \clearpage
  \fi
}{}
\newcommand{\Sigla}[2]{\textbf{#1} & #2 \\[4pt]}

% Comandos de Seção do Corpo
\newcommand{\Secao}[1]{%
  \ifSecoesEmNovaPagina\clearpage\fi
  \section{#1}%
}
\let\secaoprincipal\Secao % Compatibilidade com versão anterior
\newcommand{\Subsecao}[1]{\subsection{#1}}
\newcommand{\Subsubsecao}[1]{\subsubsection{#1}}
\newcommand{\SecaoSemNumero}[1]{%
  \ifSecoesEmNovaPagina\clearpage\fi
  \phantomsection
  \section*{#1}%
  \addcontentsline{toc}{section}{#1}%
}

% Ambientes de Listas Padronizadas (com recuo de 1,25 cm e sem espaçamento extra)
\newenvironment{ListaAlfab}{%
  \begin{enumerate}[label=\alph*),leftmargin=1.25cm,labelsep=.5em,itemsep=0pt,parsep=0pt,topsep=0pt]%
}{%
  \end{enumerate}%
}
\newenvironment{ListaNumerada}{%
  \begin{enumerate}[label=\arabic*.,leftmargin=1.25cm,labelsep=.5em,itemsep=0pt,parsep=0pt,topsep=0pt]%
}{%
  \end{enumerate}%
}
\newenvironment{ListaMarcadores}{%
  \begin{itemize}[leftmargin=1.25cm,labelsep=.5em,itemsep=0pt,parsep=0pt,topsep=0pt]%
}{%
  \end{itemize}%
}

% Ambiente para Citação Direta Longa (mais de 3 linhas: recuo de 4 cm, fonte 10, simples)
\newenvironment{CitacaoLonga}{%
  \par\vspace{\baselineskip}%
  \begingroup
  \leftskip=4cm\relax
  \singlespacing
  \fontsize{10}{12}\selectfont
  \setlength{\parindent}{0pt}%
  \noindent\ignorespaces
}{%
  \par\endgroup
  \vspace{\baselineskip}\noindent\ignorespacesafterend
}

% Ambiente para QUADRO (fechado com grade, para textos/dados qualitativos/cronograma)
% Uso: \begin{Quadro}{Título do quadro}{Fonte}{qua:identificador} ... \end{Quadro}
\NewDocumentEnvironment{Quadro}{m m m +b}{%
  \begin{center}
  \begin{minipage}{\linewidth}
    \singlespacing
    \refstepcounter{quadro}\label{#3}%
    \addcontentsline{qua}{quadro}{\protect\numberline{\thequadro}#1}%
    Quadro \thequadro{} \textendash{} #1\par
    \vspace{6pt}
    {\small
    \renewcommand{\arraystretch}{1.3}%
    \renewcommand{\tabularxcolumn}[1]{m{##1}}%
    #4
    }\par\vspace{6pt}
    {\footnotesize Fonte: #2.\par}
  \end{minipage}
  \end{center}
}{}

% Ambiente para TABELA padrão IBGE/ABNT (aberta nas laterais, para dados numéricos)
% Uso: \begin{Tabela}{Título da tabela}{Fonte}{tab:identificador} ... \end{Tabela}
\NewDocumentEnvironment{Tabela}{m m m +b}{%
  \begin{center}
  \begin{minipage}{\linewidth}
    \singlespacing
    \refstepcounter{table}\label{#3}%
    \addcontentsline{lot}{table}{\protect\numberline{\thetable}#1}%
    Tabela \thetable{} \textendash{} #1\par
    \vspace{6pt}
    {\small
    \renewcommand{\arraystretch}{1.3}%
    #4
    }\par\vspace{6pt}
    {\footnotesize Fonte: #2.\par}
  \end{minipage}
  \end{center}
}{}

% Comando para FIGURA padrão ABNT (título em cima, imagem no centro, fonte embaixo)
% Uso: \InserirFigura{Título}{caminho/arquivo.png}{0.7\linewidth}{Fonte}{fig:id}
\newcommand{\InserirFigura}[5]{%
  \begin{center}
  \begin{minipage}{\linewidth}
    \singlespacing\centering
    \refstepcounter{figure}\label{#5}%
    \addcontentsline{lof}{figure}{\protect\numberline{\thefigure}#1}%
    Figura \thefigure{} \textendash{} #1\par
    \vspace{6pt}
    \IfFileExists{#2}{%
      \includegraphics[width=#3]{#2}\par
    }{%
      \fbox{\parbox{#3}{\centering\footnotesize\vspace{12pt}%
        [Imagem não encontrada: \textbf{#2}]\\[4pt]%
        Coloque o arquivo de imagem nessa pasta ou ajuste o caminho no comando \textbf{\textbackslash InserirFigura}.\vspace{12pt}}}\par
    }%
    \vspace{6pt}
    {\footnotesize Fonte: #4.\par}
  \end{minipage}
  \end{center}
}

% Comandos para Referências, Apêndices e Anexos
\newcommand{\referencia}[1]{%
  \begin{samepage}\noindent #1\par\end{samepage}\par\addvspace{\baselineskip}%
}
\newcommand{\Apendice}[2]{%
  \clearpage
  \phantomsection
  \section*{APÊNDICE #1 \textendash{} #2}%
  \addcontentsline{toc}{section}{APÊNDICE #1 \textendash{} #2}%
}
\newcommand{\Anexo}[2]{%
  \clearpage
  \phantomsection
  \section*{ANEXO #1 \textendash{} #2}%
  \addcontentsline{toc}{section}{ANEXO #1 \textendash{} #2}%
}


\begin{document}

% ==============================================================================
% PARTE 3 — PÁGINAS INICIAIS (PRÉ-TEXTUAIS)
% ==============================================================================
% A linha abaixo monta automaticamente a Capa e a Folha de Rosto com os dados
% que você preencheu na PARTE 1 (respeitando o {sim}/{nao} da PARTE 2).
\GerarCapaEFolhaDeRosto

% ------------------------------------------------------------------------------
% DEDICATÓRIA (Opcional — para ativar, mude \MostrarDedicatoria{sim} na PARTE 2)
% ------------------------------------------------------------------------------
\begin{BlocoDedicatoria}
[TEXTO DA DEDICATÓRIA — Exemplo: Dedico este trabalho à minha família e a todos que apoiaram minha jornada acadêmica.]
\end{BlocoDedicatoria}

% ------------------------------------------------------------------------------
% AGRADECIMENTOS (Opcional — para ativar, mude \MostrarAgradecimentos{sim} na PARTE 2)
% ------------------------------------------------------------------------------
\begin{BlocoAgradecimentos}
[TEXTO DOS AGRADECIMENTOS — Escreva aqui seus agradecimentos à instituição, orientação, colegas, família e agências de fomento, se houver.]
\end{BlocoAgradecimentos}

% ------------------------------------------------------------------------------
% EPÍGRAFE (Opcional — para ativar, mude \MostrarEpigrafe{sim} na PARTE 2)
% ------------------------------------------------------------------------------
\begin{BlocoEpigrafe}{[AUTOR DA FRASE, ANO]}
[TEXTO DA EPÍGRAFE / CITAÇÃO INICIAL]
\end{BlocoEpigrafe}

% ------------------------------------------------------------------------------
% RESUMO EM PORTUGUÊS (Controlado por \MostrarResumo{sim} na PARTE 2)
% ------------------------------------------------------------------------------
% Coloque as palavras-chave no primeiro par de chaves { } separadas por ponto e
% vírgula. Escreva o corpo do resumo em um único parágrafo dentro do bloco.
\begin{BlocoResumo}{[PALAVRA 1]; [PALAVRA 2]; [PALAVRA 3].}
[TEXTO DO RESUMO — Apresente de forma concisa o contexto, o objetivo, a metodologia utilizada e os resultados obtidos ou esperados, em um único parágrafo.]
\end{BlocoResumo}

% ------------------------------------------------------------------------------
% ABSTRACT / LÍNGUA ESTRANGEIRA (Opcional — mude \MostrarAbstract{sim} na PARTE 2)
% ------------------------------------------------------------------------------
\begin{BlocoAbstract}{[KEYWORD 1]; [KEYWORD 2]; [KEYWORD 3].}
[ABSTRACT TEXT — English translation of the abstract in a single paragraph.]
\end{BlocoAbstract}

% ------------------------------------------------------------------------------
% LISTA DE ABREVIATURAS E SIGLAS (Opcional — mude \MostrarListaSiglas{sim} na PARTE 2)
% ------------------------------------------------------------------------------
% Adicione uma linha \Sigla{SIGLA}{Significado por extenso} para cada item,
% mantendo em ordem alfabética.
\begin{BlocoSiglas}
  \Sigla{[SIGLA 1]}{[Significado completo da sigla 1]}
  \Sigla{[SIGLA 2]}{[Significado completo da sigla 2]}
\end{BlocoSiglas}

% A linha abaixo gera automaticamente as Listas (Figuras, Quadros, Tabelas) e o
% Sumário conforme suas escolhas na PARTE 2, e inicia a paginação do corpo:
\GerarListasESumario


% ==============================================================================
% PARTE 4 — CORPO DO TRABALHO (SEÇÕES TEXTUAIS)
% ==============================================================================
% COMO EDITAR E REORGANIZAR AS SEÇÕES:
% - Cada seção principal começa com \Secao{TÍTULO}.
% - Subdivisões usam \Subsecao{Título} (ex.: 1.1) e \Subsubsecao{Título} (ex.: 1.1.1).
% - NUNCA digite os números nos títulos: o LaTeX calcula a sequência e atualiza
%   o Sumário sozinho!
% - Você pode renomear, mudar de lugar, duplicar ou apagar qualquer seção inteira
%   (do comentário INÍCIO até o comentário FIM correspondente).


% ==============================================================================
% INÍCIO DA SEÇÃO: PROBLEMA DE PESQUISA (ou INTRODUÇÃO)
% ==============================================================================
\Secao{PROBLEMA DE PESQUISA}
% Escreva aqui a contextualização do tema, a delimitação e a pergunta central
% que o seu trabalho busca responder. Lembre-se: deixe uma linha em branco para
% iniciar um novo parágrafo.

[TEXTO DO PROBLEMA DE PESQUISA]

% FIM DA SEÇÃO: PROBLEMA DE PESQUISA
% ==============================================================================


% ==============================================================================
% INÍCIO DA SEÇÃO: JUSTIFICATIVA
% ==============================================================================
\Secao{JUSTIFICATIVA}
% Explique a relevância científica, acadêmica, social ou regional da pesquisa
% e por que é importante realizá-la.

[TEXTO DA JUSTIFICATIVA]

% FIM DA SEÇÃO: JUSTIFICATIVA
% ==============================================================================


% ==============================================================================
% INÍCIO DA SEÇÃO: OBJETIVOS
% ==============================================================================
\Secao{OBJETIVOS}

\Subsecao{Objetivo geral}
[TEXTO DO OBJETIVO GERAL — Comece com um verbo no infinitivo, ex.: Investigar, Analisar, Avaliar, Desenvolver...]

\Subsecao{Objetivos específicos}
% Para adicionar mais itens, basta criar uma nova linha com \item.
% Para remover, apague a linha \item correspondente.
\begin{ListaAlfab}
  \item {}[OBJETIVO ESPECÍFICO 1];
  \item {}[OBJETIVO ESPECÍFICO 2];
  \item {}[OBJETIVO ESPECÍFICO 3].
\end{ListaAlfab}

% FIM DA SEÇÃO: OBJETIVOS
% ==============================================================================


% ==============================================================================
% INÍCIO DA SEÇÃO: HIPÓTESES E METODOLOGIA
% ==============================================================================
\Secao{HIPÓTESES E METODOLOGIA}
% DICA: Se o seu trabalho não tiver hipóteses, basta apagar o bloco \Subsecao{Hipóteses}
% abaixo e mudar o título acima para \Secao{METODOLOGIA}. Você também pode
% incluir uma seção de \Secao{REFERENCIAL TEÓRICO} antes da metodologia!

\Subsecao{Hipóteses}
[TEXTO DAS HIPÓTESES]

\Subsecao{Metodologia}
% Os tópicos abaixo são sugestões. Renomeie, apague ou adicione novos comandos
% \Subsubsecao{...} de acordo com a natureza da sua pesquisa.

\Subsubsecao{Caracterização da pesquisa}
[TEXTO DA CARACTERIZAÇÃO DA PESQUISA]

\Subsubsecao{Materiais, fontes ou participantes}
[TEXTO DOS MATERIAIS, FONTES OU PARTICIPANTES]

\Subsubsecao{Procedimentos}
[TEXTO DOS PROCEDIMENTOS]

\Subsubsecao{Análise dos dados}
[TEXTO DA ANÁLISE DOS DADOS]

% FIM DA SEÇÃO: HIPÓTESES E METODOLOGIA
% ==============================================================================


% ==============================================================================
% INÍCIO DA SEÇÃO: RESULTADOS ESPERADOS (ou RESULTADOS E DISCUSSÃO)
% ==============================================================================
\Secao{RESULTADOS ESPERADOS}
[TEXTO DOS RESULTADOS ESPERADOS]

% FIM DA SEÇÃO: RESULTADOS ESPERADOS
% ==============================================================================


% ==============================================================================
% INÍCIO DA SEÇÃO: CRONOGRAMA (Controlado por \MostrarCronograma{sim} na PARTE 2)
% ==============================================================================
\ifMostrarCronograma
\Secao{CRONOGRAMA}
[TEXTO DE APRESENTAÇÃO DO CRONOGRAMA — Exemplo: As etapas previstas para o desenvolvimento da pesquisa estão distribuídas no Quadro~\ref{qua:cronograma}.]

% COMO EDITAR O QUADRO DO CRONOGRAMA:
% - O sinal & separa as colunas; o comando \\ \hline encerra a linha.
% - Marque com X os períodos em que cada atividade será realizada e deixe em
%   branco (ou coloque um traço –) nos demais períodos.
% - Para ADICIONAR UMA ATIVIDADE: copie uma linha inteira antes de \end{tabularx}.
% - Para MUDAR O NÚMERO DE PERÍODOS (ex.: de 4 para 6 meses):
%   1) Na linha \begin{tabularx}, troque *{4}{M{1.4cm}|} por *{6}{M{0.9cm}|}.
%   2) Adicione mais duas colunas (& \textbf{[P5]} & \textbf{[P6]}) no cabeçalho
%      e (& & ) em TODAS as linhas de atividades.
\begin{Quadro}{[TÍTULO DO CRONOGRAMA]}{[FONTE DO CRONOGRAMA]}{qua:cronograma}
  \begin{tabularx}{\linewidth}{|L|*{4}{M{1.4cm}|}}
    \hline
    \textbf{Atividades} & \textbf{[P1]} & \textbf{[P2]} & \textbf{[P3]} & \textbf{[P4]} \\ \hline
    [ATIVIDADE 1] & & & & \\ \hline
    [ATIVIDADE 2] & & & & \\ \hline
    [ATIVIDADE 3] & & & & \\ \hline
  \end{tabularx}
\end{Quadro}
\fi
% FIM DA SEÇÃO: CRONOGRAMA
% ==============================================================================


% ==============================================================================
% PARTE 5 — ELEMENTOS PÓS-TEXTUAIS (REFERÊNCIAS, APÊNDICES E ANEXOS)
% ==============================================================================

% ------------------------------------------------------------------------------
% REFERÊNCIAS (Controladas por \MostrarReferencias{sim} na PARTE 2)
% ------------------------------------------------------------------------------
% COMO PREENCHER:
% - Coloque cada obra dentro de um comando \referencia{...} em ordem alfabética.
% - O espaçamento simples, o alinhamento à esquerda e o espaço entre uma obra e
%   outra já são aplicados automaticamente!
% - Para links da internet, use \url{https://...} dentro da referência.
% - Na PARTE 6 (no final deste arquivo), há modelos prontos de como formatar
%   Artigos, Livros, Capítulos e Leis segundo a ABNT.
\ifMostrarReferencias
\SecaoSemNumero{REFERÊNCIAS}
\begingroup
\singlespacing
\setlength{\parindent}{0pt}
\setlength{\parskip}{0pt}
\raggedright

\referencia{[REFERÊNCIA COMPLETA 1]}

\referencia{[REFERÊNCIA COMPLETA 2]}

\endgroup
\fi

% ------------------------------------------------------------------------------
% APÊNDICES (Opcional — mude \MostrarApendices{sim} na PARTE 2 para exibir)
% ------------------------------------------------------------------------------
% Use para documentos ou materiais elaborados por VOCÊ (ex.: questionários,
% roteiros, tabelas detalhadas extras). Para adicionar outro, use \Apendice{B}{...}.
\ifMostrarApendices
\Apendice{A}{[TÍTULO DO APÊNDICE A]}
[CONTEÚDO DO APÊNDICE A]
\fi

% ------------------------------------------------------------------------------
% ANEXOS (Opcional — mude \MostrarAnexos{sim} na PARTE 2 para exibir)
% ------------------------------------------------------------------------------
% Use para documentos não elaborados por você (ex.: parecer do comitê de ética,
% mapas oficiais, laudos externos). Para adicionar outro, use \Anexo{B}{...}.
\ifMostrarAnexos
\Anexo{A}{[TÍTULO DO ANEXO A]}
[CONTEÚDO DO ANEXO A]
\fi


% ==============================================================================
% PARTE 6 — CAIXA DE FERRAMENTAS (RECEITAS PRONTAS PARA COPIAR E COLAR)
% ==============================================================================
% Tudo aqui embaixo está comentado com % e NÃO aparece no PDF.
% Quando precisar de algum elemento, copie o bloco desejado, cole na seção do
% seu texto onde quer que ele apareça e apague o símbolo % do início das linhas!
%
% ------------------------------------------------------------------------------
% 1. NOVA SEÇÃO, SUBSEÇÃO E SUBSUBSEÇÃO
% ------------------------------------------------------------------------------
% \Secao{[TÍTULO DA NOVA SEÇÃO]}
% [Texto da seção...]
%
% \Subsecao{[Título da subseção]}
% [Texto da subseção...]
%
% \Subsubsecao{[Título da subsubseção]}
% [Texto da subsubseção...]
%
% ------------------------------------------------------------------------------
% 2. DIFERENÇA ENTRE QUADRO E TABELA (PADRÃO ABNT / IBGE)
% ------------------------------------------------------------------------------
% - Use QUADRO quando o conteúdo for predominantemente TEXTUAL ou qualitativo
%   (possui linhas verticais fechando as laterais e separando colunas).
% - Use TABELA quando o conteúdo for numérico/estatístico (padrão IBGE: as
%   laterais são ABERTAS, sem linhas verticais, apenas com filetes horizontais
%   no topo, abaixo do cabeçalho e na base).
% - Tipos de colunas disponíveis:
%     L = coluna elástica alinhada à esquerda (quebra linha sozinha)
%     C = coluna elástica centralizada
%     R = coluna elástica alinhada à direita
%     P{4cm} = coluna de largura fixa (4cm) alinhada à esquerda
%     M{2cm} = coluna de largura fixa (2cm) centralizada
%
% --- RECEITA DE QUADRO (Fechado com grade) ---
% \begin{Quadro}{[Título do quadro]}{[Fonte, ex.: elaboração própria (2026)]}{qua:exemplo}
%   \begin{tabularx}{\linewidth}{|P{4cm}|L|}
%     \hline
%     \textbf{[Coluna 1]} & \textbf{[Coluna 2]} \\ \hline
%     [Item A] & [Descrição detalhada do item A] \\ \hline
%     [Item B] & [Descrição detalhada do item B] \\ \hline
%   \end{tabularx}
% \end{Quadro}
% % Como citar no texto: conforme apresentado no Quadro~\ref{qua:exemplo}.
%
% --- RECEITA DE TABELA PADRÃO IBGE/ABNT (Aberta nas laterais, para números) ---
% \begin{Tabela}{[Título da tabela]}{[Fonte, ex.: dados da pesquisa (2026)]}{tab:exemplo}
%   \begin{tabularx}{\linewidth}{L c c c}
%     \hline
%     \textbf{[Amostra / Grupo]} & \textbf{[Parâmetro 1]} & \textbf{[Parâmetro 2]} & \textbf{[Média \(\pm\) DP]} \\ \hline
%     [Amostra A] & 10,5 & 2,3 & 10,5 \(\pm\) 0,2 \\
%     [Amostra B] & 12,1 & 2,8 & 12,1 \(\pm\) 0,4 \\
%     [Amostra C] & 9,8  & 2,1 & 9,8 \(\pm\) 0,3  \\ \hline
%   \end{tabularx}
% \end{Tabela}
% % Como citar no texto: de acordo com os resultados da Tabela~\ref{tab:exemplo}.
%
% ------------------------------------------------------------------------------
% 3. FIGURA PADRÃO ABNT (Título na parte superior e Fonte na parte inferior)
% ------------------------------------------------------------------------------
% % Dica: se você ainda não tiver o arquivo da imagem na pasta, o modelo não
% % dará erro; ele mostrará um quadro indicando onde colocar a imagem!
% \InserirFigura
%   {[Título da figura]}
%   {imagens/minha-figura.png}
%   {0.7\linewidth}
%   {[Fonte, ex.: elaboração própria (2026)]}
%   {fig:exemplo}
% % Como citar no texto: observa-se na Figura~\ref{fig:exemplo}.
%
% ------------------------------------------------------------------------------
% 4. EQUAÇÃO MATEMÁTICA NUMERADA + EXPLICAÇÃO DAS VARIÁVEIS
% ------------------------------------------------------------------------------
% \begin{equation}\label{eq:exemplo}
%   \mathrm{Rendimento}\;(\%) = \frac{m_{\text{final}}}{m_{\text{inicial}}} \times 100
% \end{equation}
% Em que \(m_{\text{final}}\) corresponde à massa final obtida, em gramas; e
% \(m_{\text{inicial}}\) representa a massa inicial da amostra, em gramas.
% % Como citar no texto: calculado conforme a Equação~\eqref{eq:exemplo}.
%
% ------------------------------------------------------------------------------
% 5. CITAÇÕES NO TEXTO (CURTA E LONGA — PADRÃO ABNT)
% ------------------------------------------------------------------------------
% % CITAÇÃO INDIRETA (paráfrase):
% %   Segundo Silva (2024), os resultados demonstram...
% %   ...conforme observado em estudos recentes (Silva; Santos, 2024).
% %   ...quando há mais de três autores, usa-se et al. em itálico (Sousa \textit{et al.}, 2020).
%
% % CITAÇÃO DIRETA CURTA (até 3 linhas — vai no próprio parágrafo entre aspas):
% %   Segundo Autor (2024, p. 15), ``texto exato copiado da obra original''.
%
% % CITAÇÃO DIRETA LONGA (mais de 3 linhas — recuo de 4 cm, fonte 10, simples):
% \begin{CitacaoLonga}
%   [Copie aqui o trecho literal com mais de três linhas, sem usar aspas. O recuo
%   de 4 cm à esquerda, o espaçamento simples e a fonte tamanho 10 são aplicados
%   automaticamente por este bloco] (Autor, Ano, p. [XX]).
% \end{CitacaoLonga}
%
% ------------------------------------------------------------------------------
% 6. LISTAS DE ITENS (ALFABÉTICA, NUMERADA E COM MARCADORES)
% ------------------------------------------------------------------------------
% % Lista com alíneas a), b), c):
% \begin{ListaAlfab}
%   \item {}[Primeiro item];
%   \item {}[Segundo item].
% \end{ListaAlfab}
%
% % Lista numerada 1., 2., 3.:
% \begin{ListaNumerada}
%   \item {}[Primeiro passo];
%   \item {}[Segundo passo].
% \end{ListaNumerada}
%
% % Lista com marcadores (bolinhas):
% \begin{ListaMarcadores}
%   \item {}[Ponto importante 1];
%   \item {}[Ponto importante 2].
% \end{ListaMarcadores}
%
% ------------------------------------------------------------------------------
% 7. MODELOS DE REFERÊNCIAS BIBLIOGRÁFICAS (PADRÃO ABNT NBR 6023)
% ------------------------------------------------------------------------------
% % ARTIGO DE PERIÓDICO (REVISTA CIENTÍFICA):
% \referencia{SOBRENOME, Nome; SOBRENOME, Nome. Título do artigo. \textbf{Nome da Revista em Negrito}, Cidade, v. 10, n. 2, p. 15-30, 2025. DOI: \url{https://doi.org/10.xxxx/xxxx}. Disponível em: \url{https://...}. Acesso em: 15 nov. 2026.}
%
% % ARTIGO COM MAIS DE 3 AUTORES (usa-se \textit{et al.} após o primeiro):
% \referencia{SOBRENOME, Nome \textit{et al.} Título do artigo. \textbf{Nome da Revista em Negrito}, [S. l.], v. 20, e12345, 2024. DOI: \url{https://doi.org/10.xxxx/xxxx}.}
%
% % LIVRO INTEIRO:
% \referencia{SOBRENOME, Nome. \textbf{Título do livro em negrito}: subtítulo sem negrito. 3. ed. Cidade: Editora, 2023.}
%
% % CAPÍTULO DE LIVRO ORGANIZADO POR OUTROS AUTORES:
% \referencia{SOBRENOME, Nome do Autor do Capítulo. Título do capítulo. In: SOBRENOME, Nome do Organizador (org.). \textbf{Título do livro em negrito}. Cidade: Editora, 2024. cap. 4, p. 80-102.}
%
% % LEI, INSTRUÇÃO NORMATIVA OU DOCUMENTO OFICIAL:
% \referencia{BRASIL. Ministério [Nome]. Instrução Normativa nº [XX], de [dia] de [mês] de [ano]. [Ementa/descrição]. \textbf{Diário Oficial da União}: seção 1, Brasília, DF, p. [X], [data]. Disponível em: \url{https://...}. Acesso em: 15 nov. 2026.}
%
% % TCC, DISSERTAÇÃO OU TESE:
% \referencia{SOBRENOME, Nome. \textbf{Título do trabalho em negrito}: subtítulo. 2025. Dissertação (Mestrado em [Área]) – [Faculdade/Instituto], [Universidade], [Cidade], 2025.}
% ==============================================================================

\end{document}
